\scanpage{11}
1. Morphismes représentables. Morphismes $P$-représentables.

Abus de langage : imm., im. ouverte, imm. fermée, lisse, net, étale, \unsure{Osons} dire : repr. par ....\} Critère de représentabilité par ``recollement de morceaux''. (SGAD \uline{XI} 3.5.)

2. [Faisceaux qu. \unsure{cpct|cpt}, quasi-séparés. Morphismes qu. \unsure{cpct|cpt}, qu. sép. de faisceaux]. Morphismes \inserted{loc.} de prés. finie, de prés. finie. \margin{$[\,]=$ petits caractères}

3. Morphismes $\mathrm{form}^{\mathrm{t.}}$ \unsure{lisse|lisses}, $\mathrm{form}^{\mathrm{t.}}$ nets, $\mathrm{form}^{\mathrm{t.}}$ étales, et variantes infinitésimales. Critères standart....

4. Critères pour qu'un morphisme soit un monomorphisme, un bimorphisme, une immersion ouverte.

5. Conditions de séparation:
\[
\left\{\begin{aligned}
&\text{morph. diag. repr.}\\
&\text{morphisme diagonal repr. \inserted{par imm., par imm. qu. \unsure{cpcte|cpte}, \inserted{par imm. de prés. finie,}} par imm. fermée}\\
&\text{Relations avec \uline{critères valuatifs}}
\end{aligned}\right.
\]

% The arrow brings item 6 up from the space between 7 and 8.
\inserted{6. Conditions de propreté.}

7. Espaces annelés et foncteurs \unsure{, en|sur} \unsure{``Sch''|$(\mathrm{Sch})_{/}$}. Morphismes surjectifs, \deleted{\unsure{ensemblistes}}\deleted{radiciels} radiciels (morphisme diagonal est surjectif).

[Une immersion ouverte surjective est un isom].

8. Conditions de connexion.

Principe terminologiques pour les morphismes de faisceaux sur $(\mathrm{Sch})$ : a) Conditions locales b) Si $f:F\to G$ est représentable, cela doit cadrer avec la définition déjà acceptée pour \deleted{\unsure{sché}} les \unsure{morphismes} de schémas [$\forall\ X\to G$, le morphisme de schémas $\deleted{\unsure{F}}X\times_G F\to X$ a la propriété]

% archive sheet number: 10
